\scanpage{61}
(compatible avec les op. du groupe
$G_n=\deleted{\unsure{\mathrm{Sl}(2,|\mathrm{Sl}(2,\mathbb Z)}}\mathrm{Sl}(E)/{\pm1}$, mais pas du groupe
$\mathrm{Gl}(E)$ tout entier) de $S$ dans ledit cône.

On ne trouve pas tous les éléments non dégénérés
dudit cône, mais \inserted{sur chaque rayon} seulement \deleted{\unsure{tous}} une
classe \deleted{\unsure{en}}\inserted{mod} $\mathbb F_p^*/\mathbb F_p^{\unsure{\times2|*2}}$. Les opérations \inserted{$T_\lambda\ (\lambda\in\mathbb F_p^\times)$} de
$\mathbb F_p^\times=\mathbb Z/n\mathbb Z^*$ sur $S$, sont transformées en
les opérations d'homothétie par les \uuline{$\lambda^2$}.

Mais les homothéties par les $\lambda\in\mathbb F_p^*\smallsetminus F_p^{\times2}$
échangent entre elles les deux
composantes naturelles de l'ens des \deleted{pts}
\inserted{él.} non \deleted{dégénérés}\inserted{nuls} du cône (cette division
en classes, \deleted{indexées par $k$} l'ens des
classes formant un torseur sous $k^\times/k^{\times2}$,
a un sens pour \unsure{tte|les|tt} Alg. d'Azumaya
de rang $4$ sur un corps quelconque $k$ --
si $k$ est fini, \deleted{\unsure{$q/\mathbb F_p$|$\mathbb F_q$}}, de car $\ne2$, $k^\times\deleted{\unsure{/k|/k^\times|/k^{*2}}}$
cyclique d'ordre pair, $k^\times/k^{\times2}$ est d'ordre $2$
\ldots ) Ainsi on trouve une ``bi-structure'',
avec opération de $\unsure{\mathrm{Gp}|\mathrm{GP}|\mathrm{Gl}}(E)\times\mathbb F_p^*$ dessus\ldots.

\hfill TSVP
% archive sheet number: 60
